\documentclass[11pt]{article}
\usepackage{kept-report}
\title{Business case}
\subtitle{What Kept is worth to Maple Bank, measured on the bank's own data}
\reportnumber{Business value · cost · ROI}
\runninghead{Business case}
\date{4 October 2026}

\begin{document}
\makereporttitle

\statrow{\stat{\$2.2M}{a year, base case (\$1.0M–4.2M)}\stat{3.8×}{year-one ROI, base case}\stat{\textasciitilde 2 mo}{to pay back the build}\stat{62\% vs 52\%}{promised money paid: payday-dated vs not}}

\takeaway{\textbf{The short version.} Collections loses money in one place: cases that \emph{roll} into a worse bucket.
The cheapest way to stop a roll is a promise the customer can actually keep, and the bank's own history says
promises due right after payday get \textbf{11 points more} of the promised money paid. Kept dates every promise to
payday, contacts the people a contact actually helps, swaps letters for near-free digital nudges, and gives agents
and QA reviewers time back. Counted conservatively, that's \textbf{\$2.2M a year} against \$150K to run and \$300K to
build, so it pays back in about two months. Every number below comes from the release, except six assumptions
that are named, explained and stress-tested in §7.}

\section{Where the money leaks}

\fig[0.8\linewidth]{bc_funnel.png}{Cases that entered collections at 1–30 days past due in the last 12 months, and what happened next.}

\textbf{What the chart says.} Of the 194,515 cases that entered at 1–30 days past due in the 12 months to 28 Sep 2026
(\$63.9M overdue), about half roll into a worse bucket within 30 days (48.8\% across 330,779 cases since 2016). Of the
rolled cases, \textbf{2.5\% end up charged off, at \$12,523 on average}; of the cases that \emph{don't} roll,
\textbf{none} do. Rolling is the only road to a charge-off, so everything Kept does is aimed at the first 30 days.

\section{The lever: timing the promise}

\fig[0.8\linewidth]{bc_payday.png}{Resolved promises by days between the customer's last salary deposit and the due date.}

\textbf{What the chart says.} Both panels slope the same way: the further a due date sits after payday, the worse the
promise does. Due 0–2 days after payday, \textbf{54.3\%} of promises are kept and \textbf{62.5\%} of the promised money
arrives. Due 8+ days after, when the pay is already spent, that drops to 46.3\% and \textbf{51.5\%}. The money panel
moves more than the count panel, because a late-dated promise that does get paid is more often paid only in part.

\begin{table}[htbp]\small\centering
\begin{tabular}{@{}lrrrr@{}}
\toprule
\textbf{Due date vs payday} & \textbf{Promises} & \textbf{Kept} & \textbf{Money paid} & \textbf{Share of promised \$} \\
\midrule
0–2 days after & 8,359 & 54.3\% & 62.5\% & 26\% \\
3–7 days after & 7,553 & 52.2\% & 60.9\% & 24\% \\
8+ days after & 16,045 & 46.3\% & 51.5\% & 50\% \\
\bottomrule
\end{tabular}
\caption{Where we can see the customer's salary (32,000 resolved promises).}
\end{table}

\paragraph{From the chart to dollars} Today half of the promised money sits in the worst bucket. Moving the 3+ day
promises to 0–2 days, at the 0–2 day paid rate, lifts the paid share by \textbf{5.9 points} on the \textbf{\$41.7M}
promised in the last 12 months. That's \$2.46M. We count \textbf{half} of it, \textbf{\$1.23M}, because the gap is
descriptive: part of it is \emph{who} gets offered which date, not the date itself. The 5\% hold-out measures the real
effect after launch.

\insight{How Kept uses this}{Kept sizing is plain arithmetic, not a model: the due date is 1–2 days after the expected
payroll date (detected from salary deposits), and the amount is split into up to three instalments that each fit inside
the customer's spare cash for that pay period. A promise that can't be sized that way isn't proposed at all.}

\section{A cheaper, kinder contact mix}

\fig[0.85\linewidth]{bc_mix.png}{Tonight's 96,000 open cases: the bank's current strategy vs Kept.}

\textbf{What the chart says.} Kept doesn't simply call less: call-type actions are almost unchanged (40,601 vs 41,574).
What changes is \emph{which} calls. The current strategy's biggest bar is collection calls (34,326); Kept's is
\textbf{hardship reviews (37,409)}, because text and cash-flow signals flag far more customers who need a protective
conversation than the old rule set did. At the cheap end, \$1.10 letters (7,180) disappear in favour of SMS payment links
(\$0.01), emails and app pushes (\$0.001).

\begin{table}[htbp]\small\centering
\begin{tabular}{@{}lrrr@{}}
\toprule
& \textbf{Cost per contact} & \textbf{Current strategy} & \textbf{Kept} \\
\midrule
Call-type (call, promise, plan, hardship review) & \$1.75 & 41,574 & 40,601 \\
Letter & \$1.10 & 7,180 & 0 \\
SMS payment link & \$0.01 & 27,942 & 11,971 \\
Email · app push & \$0.001 & 0 & 20,370 \\
\midrule
\textbf{Cost of one pass over the open book} & & \textbf{\$80,932} & \textbf{\$71,192} \\
\bottomrule
\end{tabular}
\caption{Channel costs are the release's own averages for September 2026.}
\end{table}

Each pass over the book saves \textbf{\$9,740}. The open book is worked through roughly weekly under the
3-contacts-in-7-days cap, so 52 passes a year give \textbf{\$506K}.

\paragraph{What the model expects} The uplift model estimates that tonight's 55,674 positive-uplift proposals bring in
\textbf{\$2.06M more within 14 days} than doing nothing. That's a model estimate, so it's \emph{not} in our totals; it's
exactly what the hold-out is there to check.

\section{Time back for agents and QA}

\begin{table}[H]\small\centering
\begin{tabularx}{\linewidth}{@{}lrL@{}}
\toprule
& \textbf{Base} & \textbf{How it's computed} \\
\midrule
Call summaries & \$378K & 629,435 agent calls a year × 60 s of after-call work saved × \$36/h loaded agent cost (≈ 5.8 agents' time back on the phones) \\
QA pre-scoring & \$45K & today's 2\% manual sample is 2,518 reviewer-hours a year; scoring every call first lets reviewers skip the clean ones (50\% saved) and coverage goes from 2\% of calls to all of them \\
\bottomrule
\end{tabularx}
\caption{Average handle time in the release is 6.8 minutes, so 60 seconds of wrap-up is a modest assumption.}
\end{table}

\section{The upside we don't count: roll rate}

From §1: 194,515 cases a year enter at 1–30 days past due, 2.5\% of the ones that roll get charged off, at \$12.5K each.
So \textbf{every point of 30-day roll rate is worth about \$603K a year} in charge-offs avoided, before counting the
collection cost of the rolled cases. Our design hypothesised 1–2 points.

We keep it \textbf{out of the base case}: a kept promise \emph{is} what stops a roll, so counting both would count the
payday effect twice. The high case adds one point.

\section{Putting it together}

\fig[0.9\linewidth]{bc_scenarios.png}{Annual value by component and scenario.}

\textbf{What the chart says.} In every scenario the payday effect is the largest block, and the contact mix the second.
The low case still clears \$1M a year with half the payday capture, half the book passes and half the time saved. The
high case more than doubles the base mostly because of the roll-rate point it adds.

\begin{table}[htbp]\small\centering
\begin{tabular}{@{}lrrr@{}}
\toprule
\textbf{Annual value (CAD)} & \textbf{Low} & \textbf{Base} & \textbf{High} \\
\midrule
Promises dated to payday & 614,603 & 1,229,207 & 1,843,810 \\
Cheaper contact mix & 253,239 & 506,479 & 1,012,957 \\
Agent wrap-up time (call summaries) & 157,359 & 377,661 & 660,907 \\
QA reviewer time & 18,883 & 45,319 & 79,309 \\
Roll-rate reduction (charge-offs avoided) & 0 & 0 & 603,362 \\
\midrule
\textbf{Total benefit a year} & \textbf{1,044,084} & \textbf{2,158,666} & \textbf{4,200,345} \\
Run cost a year & 150,000 & 150,000 & 150,000 \\
Build, one-off & 300,000 & 300,000 & 300,000 \\
\midrule
\textbf{Year-one ROI} & \textbf{1.3×} & \textbf{3.8×} & \textbf{8.3×} \\
\textbf{Payback} & 4.0 months & 1.8 months & 0.9 months \\
\bottomrule
\end{tabular}
\caption{Year-one ROI = (benefit − run − build) ÷ (run + build). Payback = build ÷ (benefit − run), in months.}
\end{table}

\paragraph{Why it's cheap to run} One nightly batch scores the whole book in about seven minutes on a single
machine. Text runs on a local open model (no per-call fees, and notes never leave the bank); a hosted model is only used
for the Ask. DuckDB and dbt replace a cluster. The \$150K is mostly one operations engineer (\$120K) plus about \$30K of
compute and hosted-model use at production volume. The \$300K build is four engineers for three months.

\section{What moves the number}

\fig[0.88\linewidth]{bc_sensitivity.png}{Each assumption swung from its low to its high value, everything else at base.}

\textbf{What the chart says.} The answer hinges on one question: \textbf{how much of the payday gap is causal?} That
assumption swings the total by about \$1.2M. The next biggest are how often the book is worked (contact savings) and
the roll-rate point. Minutes saved and QA time barely move it. That's why the rollout plan is built around measuring the
payday effect first.

\begin{table}[htbp]\small\centering
\begin{tabularx}{\linewidth}{@{}lrrrL@{}}
\toprule
\textbf{Assumption} & \textbf{Low} & \textbf{Base} & \textbf{High} & \textbf{Why this range} \\
\midrule
Payday gap we capture & 25\% & 50\% & 75\% & the gap is observational; half is a cautious read \\
Passes over the book a year & 26 & 52 & 104 & weekly work-through under the contact cap \\
After-call work saved & 30 s & 60 s & 90 s & the summary pre-fills the wrap-up note \\
Loaded agent cost & \$30/h & \$36/h & \$42/h & \textasciitilde \$50K salary + 30\% on-cost \\
QA reviewer time saved & 25\% & 50\% & 75\% & reviewers skip pre-scored clean calls \\
Roll-rate points counted & 0 & 0 & 1 & overlaps the payday effect \\
\bottomrule
\end{tabularx}
\caption{The only assumptions in the model. Everything else is measured.}
\end{table}

\section{Risks, costs we're honest about, and how we'll know}

\begin{itemize}
\item \textbf{Specialist workload.} 38,115 of tonight's cards (40\%) are Red: a hardship specialist reviews before anything
is sent. That's the protection working, but it's real time. A rule for customers already on an active hardship plan would
cut it, and it's the next thing we'd ship.
\item \textbf{Observational evidence.} The payday gap and the model's \$2.06M estimate are both observational. A permanent
\textbf{5\% hold-out} stays on the old strategy, so the real effect on promises kept, roll rate and cost is measured
from day 30, and a \textbf{kill switch} reverts everyone if it isn't there.
\item \textbf{What isn't in the numbers.} Fewer complaints and less regulator risk from finding hardship when the customer
says it rather than weeks later; fairness checked on every run; a plain-English reason and an audit trail for every
decision. Those are hard to price and matter anyway.
\end{itemize}

\reproduce{\path{scripts/business_case.py} reads only the release files and the scored cards, and caches the
numbers in \path{data/exports/reports/business_data.json}.}{\code{uv run python scripts/reportlib.py business\_case}}

\end{document}
